% ==============================================================================
% CAPÍTULO 5: CONCLUSÕES E TRABALHOS FUTUROS
% ==============================================================================

\section{CONCLUSÕES E TRABALHOS FUTUROS}
\label{sec:capitulo5_conclusoes}

Esta dissertação desenvolveu, parametrizou e avaliou \textit{in silico} um ambiente computacional integrado para o gerenciamento térmico-energético de um nó de centro de dados orbital de $50\text{ kW}$ ($100$ módulos aceleradores NVIDIA A100 SXM4 de $400\text{ W}$ de TDP nominal), operando no vácuo em órbita heliossíncrona (\textit{Sun-Synchronous Orbit} -- SSO a $700\text{ km}$, LTAN $18\text{h}00$, período orbital $T_{orb} = 5.926{,}38\text{ s}$ e eclipse máximo de $1.105{,}39\text{ s}$). A ausência de convecção atmosférica, a rejeição térmica estritamente radiativa governada pela Lei de Stefan--Boltzmann e as transições periódicas entre insolação e eclipse impõem um acoplamento rígido entre geração fotovoltaica, armazenamento eletroquímico e controle ativo de potência computacional.

\subsection{Síntese dos Resultados e Fechamento dos Objetivos}
\label{sec:sintese_objetivos_cap5}

Os resultados analíticos e numéricos consolidados nos Capítulos~3 e~4 sintetizam-se em quatro eixos principais:

\begin{enumerate}[label=\textbf{\arabic*.},leftmargin=*]
    \item \textbf{Modelagem Multifísica Integrada e Contrato Único de Parâmetros (Capítulo~3):} A plataforma (\texttt{PROJETO\_1P0}) acoplou a mecânica orbital com eclipses analíticos, o arranjo fotovoltaico em células \textbf{Spectrolab XTJ} ($29{,}5\%$ BOL, $27.073$ células, $161{,}49\text{ m}^2$ ativos / $179{,}43\text{ m}^2$ totais para $50\text{ kW}$ EOL a $65^\circ\text{C}$), a conversão estática de potência, o banco de baterias \textbf{Saft VES~16} ($19\text{P}108\text{S}$, $2.052$ células, $236{,}0\text{ kg}$, $E_{EOL} = 26{,}59\text{ kWh}$ operando em $SoC \in [0{,}70,\; 0{,}95]$), o modelo eletrotérmico DVFS com fuga térmica exponencial e a planta térmica acoplada (nó de junção, coletor/fluido e aleta radiativa 1D com revestimento Z-93 EOL, $\alpha_s = 0{,}36$, $\varepsilon = 0{,}90$, $A_{pl} = 49{,}90\text{ m}^2$ / $A_{emit} = 99{,}80\text{ m}^2$). Sem limitação da frequência de operação (\texttt{S0}), a temperatura de junção ultrapassa o limite de projeto de $80^\circ\text{C}$.

    \item \textbf{Modelo Substituto PINN e Estimador de Estado Térmico (Seção~\ref{sec:resultados_marco4}):} A rede neural informada por física (\texttt{v2c.pt}, retreinada sobre a planta fisicamente consistente com $c_a = 4.860\text{ J/(m}^2\cdot\text{K)}$) atingiu erro mediano de $0{,}215^\circ\text{C}$, $89{,}3\%$ dos casos em $\le 0{,}5^\circ\text{C}$ e pior subestimação de $-1{,}086^\circ\text{C}$ no conjunto independente no domínio ($n = 300$), com inferência isolada média de $0{,}36\text{ ms}$ ($p_{99} = 0{,}39\text{ ms}$, $252{,}2\times$ mais rápida que o solver EDP $N=160$) em CPU com uma \textit{thread}. Acoplada ao estimador por mínimos quadrados ponderados para a diferença de temperatura raiz-ponta ($\hat{\Delta}_{ponta}$, não instrumentada em voo), o erro RMS em $T_j(300\text{ s})$ foi reduzido de $1{,}195\text{ K}$ para $0{,}274\text{ K}$ sob ruído e viés de sensor ($\pm 1{,}5\text{ K}$).

    \item \textbf{Desempenho em Malha Fechada na Campanha Definitiva de 20 Órbitas e 20 Sementes (Seção~\ref{sec:resultados_marco5}):} Em \textbf{1.728 simulações} definitivas (20 órbitas, 20 sementes de teste independentes das de calibração, 10–29, com intervalos de confiança \textit{t} de Student pareados e correção de Bonferroni), o Sistema de Gerenciamento Térmico-Energético (TEMS), controlador preditivo cuja única variável manipulada é a frequência de relógio (DVFS), apresentou resultados bem delimitados frente aos cinco critérios de avaliação:
    \begin{itemize}
        \item \textbf{P1 (ATENDIDO):} zero violações térmicas nas 480 execuções TEMS (24 grupos) nos regimes nominal, de rajadas, de quantização NVML e fora de domínio (pior $T_j = 79{,}61^\circ\text{C}$) com a margem da regra de cada preditor ($m_{pred} = 1{,}09\text{ K}$ na PINN; $0{,}33\text{ K}$ no RC3), enquanto o controlador reativo PI sintonizado (\texttt{R\_SINT}) violou $80^\circ\text{C}$ em 4 grupos (até $80{,}34^\circ\text{C}$). Sem margem ($m_T = 0$), houve violações nos dois preditores; com $m_T = 0{,}5\text{ K}$, houve violações na PINN.
        \item \textbf{P2 (NÃO ATENDIDO):} sob redução isolada de emissividade (\texttt{eps\_baixo}, $\varepsilon: 0{,}90 \to 0{,}80$), \texttt{TEMS\_PINN} manteve $T_j \le 80^\circ\text{C}$ em 20/20 sementes ($78{,}97^\circ\text{C}$), contra 1/20 violação marginal do \texttt{TEMS\_RC3} ($80{,}03^\circ\text{C}$) e 20/20 violações do \texttt{R\_SINT} ($81{,}43^\circ\text{C}$). Porém, sob aumento da resistência térmica interna (\texttt{r\_int\_alto}, \texttt{r\_int\_muito\_alto}) e no cenário \texttt{combinada} ($R_{int}, C_{mod}, \varepsilon$), ambos os preditores do TEMS violaram $80^\circ\text{C}$ em 20/20 sementes por não estimarem $R_{int}$ online.
        \item \textbf{P3 (ATENDIDO):} zero violações do piso energético ($SoC_{\min} = 0{,}700$, resíduo máximo de barramento de $1{,}46 \times 10^{-11}\text{ W}$) em todas as \textbf{1.586 execuções} com supervisor ativo.
        \item \textbf{P4 (ATENDIDO COM EMPATE):} todos os quatro grupos na fronteira de Pareto ficaram em empate estatístico entre TEMS e \texttt{R\_SINT} pelo IC \textit{t} pareado com correção de Bonferroni (20 pares sem calibração; 12 pares comparáveis e 8 censurados no sensor calibrado), sem acionamento da cláusula de falsificação (nem estrita nem pós-hoc em $0{,}25\text{ K}$).
        \item \textbf{P5 (ATRIBUÍVEL EM ALGUNS GRUPOS) e Latência:} a PINN não produziu diferença de vazão atribuível em 11 de 19 grupos frente ao modelo concentrado RC3 (nos 8 restantes, vazão $-0{,}20\%$ a $-0{,}91\%$ menor, reflexo direto da margem preditiva mais conservadora $1{,}09\text{ K}$ vs. $0{,}33\text{ K}$), reduzindo em \textbf{$6{,}29\times$} o tempo médio de decisão completa em malha fechada ($7{,}75\text{ ms}$ contra $48{,}75\text{ ms}$ no RC3, em CPU de servidor \texttt{i5-8500T} ociosa).
    \end{itemize}

    \item \textbf{Convergência Orbital e Periodicidade Energética:} O estudo de convergência em 20 órbitas demonstrou que a memória da condição inicial de $SoC$ desaparece até a 4ª órbita, mas a razão não inteira entre $T_{orb}$ e $\Delta t_{ctrl} = 60\text{ s}$ gera um batimento de fase com período de $\approx 4{,}4$ órbitas ($\rho_{\text{Spearman}} = -0{,}998$). Adotando a emenda pós-hoc de janela de 9 órbitas ($\approx 2$ ciclos de batimento), a periodicidade energética ($|\Delta E|/E_{TI} < 10^{-3}$) foi confirmada em 380/380 execuções de \texttt{R\_SINT} ($100{,}0\%$), 579/580 de \texttt{TEMS\_RC3} ($99{,}83\%$), 612/620 de \texttt{TEMS\_PINN} ($98{,}71\%$), 6/6 de \texttt{TEMS\_SOLVER} ($100{,}0\%$) e 141/141 de \texttt{S1} ($100{,}0\%$).
\end{enumerate}

\subsection{Limitações Epistêmicas e Fronteiras de Validade}
\label{sec:limitacoes_epistemicas_cap5}

As conclusões acima estão circunscritas às seguintes limitações declaradas:
\begin{itemize}[leftmargin=*]
    \item \textbf{Validação exclusivamente \textit{in silico}:} todos os resultados foram obtidos por simulação numérica em CPU de servidor (\texttt{i5-8500T}), sem ensaio em câmara de termovácuo (TVAC) nem execução embarcada em computador de bordo (OBC) ou FPGA qualificado para voo.
    \item \textbf{Rastreabilidade paramétrica:} parte dos parâmetros da planta térmica e do modelo eletrotérmico constitui hipóteses de projeto declaradas ($C_{mod}$, $C_{fluid}$, $R_{int}$, lei DVFS/fuga, eficiências e limiares OOD), parâmetro geométrico herdado ($L = 0{,}075\text{ m}$; com $c_a = 4.860\text{ J/(m}^2\cdot\text{K)}$ já consistente com $\delta = 2{,}0\text{ mm}$ de Al 6061-T6) e valor de referência da liga de alumínio ($k = 167\text{ W/(m}\cdot\text{K)}$).
    \item \textbf{Estrutura do controlador e simplificações físicas:} o TEMS manipula apenas a frequência comum $f$ (sem alocação heterogênea de nós nem controle de bombeamento), utiliza $f/f_{\max}$ como proxy de vazão útil (ignorando gargalos de memória), modela o BESS por balanço energético químico sem curva eletroquímica de tensão terminal no envelope $[0{,}70,\; 0{,}95]$ — razão pela qual modos de falha (FMECA) em extremos de carga (descarga profunda próximo a $0\%$ de $SoC$ ou sobrecarga acima de $95\%/100\%$) ficam fora do escopo desta dissertação — e não estima $R_{int}$ em tempo real.
\end{itemize}

\subsection{Recomendações para Trabalhos Futuros}
\label{sec:trabalhos_futuros_cap5}

Com base diretamente nas limitações identificadas nos resultados experimentais, propõem-se cinco linhas de continuidade:
\begin{enumerate}[label=\textbf{(\alph*)},leftmargin=*]
    \item \textbf{Estimação Online de $R_{int}$ e Controle Adaptativo:} incorporar um estimador recursivo da resistência térmica interna junção-coletor ($R_{int}$) acoplado ao preditor (RC3 ou PINN parametrizada em $R_{int}$), endereçando diretamente a causa raiz das violações observadas no critério P2 (\texttt{r\_int\_alto}, \texttt{r\_int\_muito\_alto} e \texttt{combinada}).
    \item \textbf{Modelagem Eletroquímica OCV/CC-CV e FMECA de Extremos de $SoC$:} estender o modelo macroscópico do BESS para uma modelagem eletroquímica OCV/CC-CV com circuito equivalente de polarização e proteção de pré-carga em baixo $SoC$ ($0\%\text{--}20\%$) para recuperação de \textit{safe mode}, habilitando a análise FMECA de subtensão/descarga profunda e de sobrecarga acima de $95\%/100\%$.
    \item \textbf{Implementação Embarcada \textit{Hardware-in-the-Loop} (HIL):} portar a PINN quantizada e a bisseção do TEMS para um processador ou FPGA tolerante à radiação, medindo o tempo real de decisão completo e o consumo energético do próprio controlador sob restrições de OBC espacial.
    \item \textbf{Co-Design de Controle Multivariável e Escalonamento via OISL:} expandir o TEMS de variável única (DVFS) para controle conjunto de frequência, desligamento seletivo de módulos e migração de cargas de IA entre satélites da constelação via enlaces ópticos inter-satélites (OISL) durante eclipses ou anomalias térmicas.
    \item \textbf{Validação Experimental em Câmara de Termovácuo (TVAC):} construir bancada em escala reduzida com aquecedores controlados por DVFS, sensores RTD/diodo com ruído real e radiador em vácuo para calibrar experimentalmente $R_{int}$, $C_{mod}$, $c_a$ e as propriedades termo-ópticas ($\alpha_s$, $\varepsilon$) fora do regime puramente simulado.
\end{enumerate}
